%========================== CONCLUSÃO ==========================
\section{Conclusão}

% ESTRUTURA EM TRÊS CAMADAS -- é o padrão adotado pelo orientador, e
% serve para não prometer mais do que o dado sustenta:
%
%   EM TERMOS DE FATOS          o que foi construído e o que se observou
%   EM TERMOS DE INTERPRETAÇÃO  o que os dados permitem afirmar
%   EM TERMOS DE HIPÓTESE       o que fica como conjectura a confirmar
%
% No caso deste trabalho a separação é especialmente importante: tudo
% o que toca eficácia terapêutica pertence à TERCEIRA camada, nunca
% às duas primeiras.
%
% O que já se pode afirmar na camada de FATOS, independentemente da
% coleta, porque está medido:
%   - o artefato foi construído e funciona nas 25 combinações de
%     animal e nível, sem falha na verificação automatizada;
%   - executa em computador, celular e óculos a partir do mesmo
%     endereço, sem instalação;
%   - sustenta 24 exemplares animados mantendo a taxa de quadros,
%     em renderização por software.

[ESCREVER]

\subsection{Contribuições}

% Separar as duas: a de engenharia e a de pesquisa.
%
% De engenharia, já sustentadas:
%   - protótipo de exposição gradual executável em navegador, com a
%     interface construída dentro da cena tridimensional;
%   - ferramenta de verificação automatizada das 25 combinações;
%   - conjunto de decisões de projeto documentadas e justificadas
%     (escala de apresentação, comportamento por exemplar, separação
%     dos apontadores).
%
% De pesquisa: depende da coleta.
[ESCREVER]

\subsection{Trabalhos futuros}

% Candidatos já levantados ao longo do desenvolvimento:
%   - rastreamento de mãos por visão computacional, usando a câmera
%     do aparelho, para permitir aproximação com a própria mão;
%   - medida fisiológica (frequência cardíaca) com sensor de perfil
%     padrão, como cinta torácica;
%   - estudo com participantes com fobia diagnosticada, mediante
%     aprovação em Comitê de Ética em Pesquisa;
%   - acompanhamento ao longo de várias sessões, em lugar de sessão
%     única;
%   - apresentação em escala real, verificando se a perda de
%     visibilidade compromete ou não o estímulo.
[ESCREVER]
